\documentclass[11pt]{amsart}

\usepackage{amsmath,amssymb,amsthm}
\usepackage{mathtools}
\usepackage[margin=1.1in]{geometry}
\usepackage{hyperref}

\newtheorem{theorem}{Theorem}[section]
\newtheorem{lemma}[theorem]{Lemma}
\newtheorem{proposition}[theorem]{Proposition}
\newtheorem{corollary}[theorem]{Corollary}
\newtheorem{conjecture}[theorem]{Conjecture}
\newtheorem{question}[theorem]{Question}
\theoremstyle{definition}
\newtheorem{definition}[theorem]{Definition}
\newtheorem{example}[theorem]{Example}
\theoremstyle{remark}
\newtheorem{remark}[theorem]{Remark}

\newcommand{\R}{\mathbb{R}}
\newcommand{\C}{\mathbb{C}}
\newcommand{\Z}{\mathbb{Z}}
\newcommand{\Q}{\mathbb{Q}}
\newcommand{\AQ}{A_Q}
\newcommand{\AP}{A_P}
\newcommand{\gapp}{\operatorname{gap}}
\newcommand{\tr}{\operatorname{tr}}
\newcommand{\wt}{\operatorname{wt}}
\newcommand{\dz}[1]{\mathrm{d}_{\Z}(#1)}
\DeclareMathOperator{\SO}{SO}
\DeclareMathOperator{\SU}{SU}

\title[Maximum-likelihood quantum automatic complexity]{Maximum-likelihood quantum automatic complexity}

\author{Bj{\o}rn Kjos-Hanssen \\ Swarnalakshmi Lakshmanan}
\address{University of Hawai\textquoteleft i at M\=anoa}
\email{bjoern.kjos-hanssen@hawaii.edu}

\date{\today}

\begin{document}

\begin{abstract}
We introduce a maximum-likelihood variant $\AQ(w)$ of quantum automatic
complexity: the least dimension $d$ such that two unitaries
$U_0,U_1\in U(d)$, with initial state $e_1$ and a designated accepting
basis state, make $w$ strictly the most likely accepted string of its
length. This stands to the exact-acceptance complexity of
\cite{KH17} as Gill's probabilistic complexity $\AP$ \cite{Gill24}
stands to deterministic complexity. We prove that $\AQ(w)=2$ for
\emph{every} nonempty string $w$, over every finite alphabet: a single
universal pair of unitaries---lifts to $\SU(2)$ of two rotations
generating a free group---witnesses every string, with acceptance
probability exactly $1$, once the initial state is aimed at the
rotation axis of $U(w)$. In the opposite direction, the
gap-parameterized complexity $A_{Q,\delta}$ is unbounded for every
fixed $\delta>0$: cube-free strings of length $n$ require dimension
$\Omega_\delta(\log n)$, and so---in sharp contrast with the
probabilistic case, where two states suffice forever---does the string
$0^n$, because unitary dynamics is recurrent and admits no absorbing
states. Finally we locate the computational difficulty of the
underlying problems: finding the most likely string of a
measure-once quantum automaton is NP-complete, with exact rational
acceptance values and an inverse-polynomial promise gap, and verifying
a claimed $\delta$-witness is coNP-complete, while the exact value of
$A_{Q,\delta}(w)$ is computable in exponential time. All estimates
are given with explicit constants and elementary proofs, with a view
toward machine formalization.
\end{abstract}

\maketitle

\section{Introduction}

Gill \cite{Gill24} defined the probabilistic automatic complexity
$\AP(w)$ of a finite string $w$: the least number of states of a
probabilistic finite automaton for which $w$ is the unique most likely
accepted string of its length. In the companion note \cite{KHpfa} it
is shown that $\AP(w)\le3$ for every string, while the parameterized
version $A_{P,\delta}$ ($\delta>0$) is unbounded. Here we carry out
the analogous program for \emph{quantum} finite automata in the
measure-once model of Moore and Crutchfield \cite{MC00}: unitary
evolution, with a single projective measurement at the end of the
input.

Our results, in brief. At $\delta=0$ the complexity collapses
completely: $\AQ(w)=2$ for every nonempty $w$
(Theorem~\ref{thm:all}). This is strictly stronger than the classical
situation, where two \emph{states} suffice only for a thin regular
family of strings \cite[Theorem~4.1]{Gill24} and the sharp universal
bound is three \cite{KHpfa}. The proof is short: if $R_0,R_1$ are
rotations of the sphere generating a free group, then the rotations
$R_z$, $z$ ranging over the strings of length $n$, are pairwise
non-commuting, so the axis of $R_w$ is fixed by $R_w$ and by no other
$R_z$; taking that axis as the initial \emph{and} accepting state
makes $w$ the unique string accepted with probability $1$. The
``memory'' of the automaton is thus a single point on the Bloch
sphere. At fixed $\delta>0$ the complexity is unbounded
(Theorems~\ref{thm:cubefree} and \ref{thm:zeros}), and here quantum
mechanics is \emph{more} constrained than probability: classically
$A_{P,\delta}(0^n)=2$ for all $n$ via an absorbing state, but unitary
dynamics has no absorbing states, is recurrent, and forces
$A_{Q,\delta}(0^n)\to\infty$. In slogan form: superposition is memory
enough for any single string, but a quantum lock wears out. The
complexity information of the measure-once model therefore resides
entirely in the gap structure, and Section~\ref{sec:complexity} shows
that this structure is computationally rich: the string-maximization
problem for measure-once automata is NP-complete and
$\delta$-witness verification is coNP-complete.

This note is written with machine formalization in mind.
Section~\ref{sec:notation} fixes all conventions, including the
recursive definitions of the basic objects.
Section~\ref{sec:toolbox} isolates the elementary ingredients: the
Lipschitz property of the Born rule, a covering bound for spheres
proved by volume comparison, simultaneous Dirichlet approximation
proved by pigeonhole, the Lyndon--Sch\"utzenberger theorem on
commuting words, and the dictionary between $\SU(2)$ and $\SO(3)$,
each with an explicit-constant statement and a complete proof. Beyond
finite-dimensional linear algebra and two classical facts---the
Lyndon--Sch\"utzenberger theorem and the existence of two rotations
generating a free group---the note is self-contained; no compactness,
measure theory, or genericity is used.

\section{Notation and conventions}\label{sec:notation}

\subsection*{Vectors and matrices}
$\C^d$ carries the inner product
$\langle x|y\rangle=\sum_{j=1}^d\overline{x_j}\,y_j$, conjugate-linear
in the \emph{first} argument, with norm $\|x\|=\langle
x|x\rangle^{1/2}$; $e_1,\dots,e_d$ is the standard basis. By
Cauchy--Schwarz, $|\langle x|y\rangle|\le\|x\|\,\|y\|$, with equality
for unit vectors if and only if $x$ and $y$ are parallel. For a
$d\times d$ complex matrix $A$, $\|A\|$ denotes the operator norm
$\sup_{\|x\|=1}\|Ax\|$; it is submultiplicative, and
$|\langle x|Ay\rangle|\le\|A\|$ whenever $\|x\|=\|y\|=1$. $A^\dagger$
is the conjugate transpose; $U(d)$ is the group of unitary matrices
($U^\dagger U=I$) and $\SU(d)$ its subgroup of determinant one.
Unitaries are isometries: $\|Ux\|=\|x\|$. Conjugation by a unitary
preserves the operator norm. We use throughout: if $U=VDV^\dagger$
with $V$ unitary and $D=\operatorname{diag}
(e^{i\theta_1},\dots,e^{i\theta_d})$, then
$\|U^p-I\|=\|D^p-I\|=\max_j|e^{ip\theta_j}-1|$.

\subsection*{Real distances}
For $t\in\R$, $\dz{t}=\min_{k\in\Z}|t-k|$ is the distance to the
nearest integer.

\subsection*{Strings}
A \emph{binary string} of length $n\ge0$ is a function
$w\colon\{1,\dots,n\}\to\{0,1\}$, written $w=w_1w_2\cdots w_n$; the set
of such is $\{0,1\}^n$, and $\{0,1\}^*$ is the free monoid of all
binary strings under concatenation, written by juxtaposition, with
the empty string as identity. For $0\le i\le j\le n$, the
\emph{factor} $w_{(i,j]}$ is the string $w_{i+1}w_{i+2}\cdots w_j$ of
length $j-i$, and $w_{[1,m]}=w_{(0,m]}$ is a \emph{prefix}. The
\emph{reversal} is $w^R_m=w_{n+1-m}$. A nonempty string $u$ is
\emph{primitive} if $u=v^k$ implies $k=1$. A string $w$ is
\emph{cube-free} if there is no nonempty $u$ with $uuu$ a factor of
$w$.

\subsection*{Automata}
Fix $U_0,U_1\in U(d)$. Define $U(z)$ for $z\in\{0,1\}^m$ recursively
by $U(\text{empty string})=I$ and
$U(z)=U_{z_m}\,U(z_{[1,m-1]})$; explicitly,
$U(z)=U_{z_m}U_{z_{m-1}}\cdots U_{z_1}$, so that the \emph{first}
letter of $z$ acts \emph{first} (rightmost in the operator
product). Given also an accepting index $k\in\{1,\dots,d\}$, the
acceptance probability of $z\in\{0,1\}^n$ is
\[
  \rho(z)=\bigl|\langle e_k|U(z)e_1\rangle\bigr|^2\in[0,1].
\]

\begin{definition}\label{def:AQ}
Let $w\in\{0,1\}^n$, $n\ge1$. A \emph{$d$-dimensional quantum witness}
for $w$ is a triple $(U_0,U_1,k)$ as above such that
\[
  \gapp(w):=\min_{z\in\{0,1\}^n\setminus\{w\}}
  \bigl(\rho(w)-\rho(z)\bigr)>0 .
\]
The \emph{maximum-likelihood quantum automatic complexity} $\AQ(w)$ is
the least $d$ admitting a $d$-dimensional quantum witness for $w$, and
for $\delta>0$,
\[
  A_{Q,\delta}(w)=\min\{\,d:\text{some $d$-dimensional witness for $w$
    has }\gapp(w)\ge\delta\,\}.
\]
\end{definition}

Thus the automaton reads $w$ left to right, is started in the basis
state $e_1$, and accepts by measuring onto the basis state $e_k$; the
case $k=1$, in which the accepting state \emph{is} the initial state,
is allowed. The exact-acceptance complexities $Q$, $Q_s$, $Q_f$ of
\cite{KH17} require the stronger condition $U(w)e_1\in\C e_k$ together
with the exclusion of all other strings of the same length from the
ray $\C e_k$; maximum-likelihood acceptance is the natural relaxation,
and the direct quantum analogue of Gill's $\AP$.

\section{Basic properties}\label{sec:basic}

\begin{proposition}\label{prop:basic}
For every nonempty binary string $w$:
\begin{enumerate}
\item $\AQ(w)\ge2$;
\item $\AQ(\bar w)=\AQ(w)$, where $\bar w$ exchanges the letters
  $0\leftrightarrow1$;
\item $\AQ(w^R)=\AQ(w)$;
\item $\AQ(w)\le Q_s(w)$, the semi-classical exact-acceptance
complexity of \cite{KH17}.
\end{enumerate}
The same statements hold for $A_{Q,\delta}$ in place of $\AQ$, for
each fixed $\delta$, except that in (4) the conclusion reads
$A_{Q,\delta_0}(w)\le Q_s(w)$ where
$\delta_0=\gapp$ of the exact witness.
\end{proposition}

\begin{proof}
(1) If $d=1$ then $U_0=(e^{i\beta_0})$ and $U_1=(e^{i\beta_1})$ are
unimodular scalars, so $U(z)$ is unimodular and $\rho(z)=1$ for every
$z$; since $n\ge1$ there are at least two strings of length $n$, and
no string is a strict maximizer.

(2) If $(U_0,U_1,k)$ witnesses $w$, then $(U_1,U_0,k)$ witnesses
$\bar w$ with the same gap, since the roles of the letters are
exchanged.

(3) Let $(U_0,U_1,k)$ witness $w$, and let $P$ be the permutation
matrix with $Pe_1=e_k$, $Pe_k=e_1$, and $Pe_j=e_j$ otherwise ($P=I$ if
$k=1$); note $P=P^\dagger=P^{-1}$ and $P$ is unitary. Put
$W_b=PU_b^\dagger P$ for $b\in\{0,1\}$; these are unitary. For any
$z\in\{0,1\}^n$, using $|\langle x|y\rangle|=|\langle y|x\rangle|$ and
inserting $P^2=I$ between all factors,
\begin{align*}
  \bigl|\langle e_k|U(z)e_1\rangle\bigr|
  &=\bigl|\langle e_1|U(z)^\dagger e_k\rangle\bigr|
   =\bigl|\langle e_1|U_{z_1}^\dagger U_{z_2}^\dagger\cdots
     U_{z_n}^\dagger e_k\rangle\bigr|\\
  &=\bigl|\langle Pe_1|\,(PU_{z_1}^\dagger P)(PU_{z_2}^\dagger P)\cdots
     (PU_{z_n}^\dagger P)\,Pe_k\rangle\bigr|
   =\bigl|\langle e_k|W_{z_1}W_{z_2}\cdots W_{z_n}\,e_1\rangle\bigr| .
\end{align*}
Now $W_{z_1}W_{z_2}\cdots W_{z_n}=W(z^R)$: indeed the recursive
definition gives
$W(z^R)=W_{(z^R)_n}\cdots W_{(z^R)_1}$ and $(z^R)_m=z_{n+1-m}$, so the
factors are $W_{z_1},W_{z_2},\dots,W_{z_n}$ from left to right. Hence
the triple $(W_0,W_1,k)$ assigns to each $z^R$ the same acceptance
probability that $(U_0,U_1,k)$ assigns to $z$; since $z\mapsto z^R$ is
a bijection of $\{0,1\}^n$, the gap of $(W_0,W_1,k)$ at $w^R$ equals
the gap of $(U_0,U_1,k)$ at $w$. This proves $\AQ(w^R)\le\AQ(w)$, and
equality follows by applying the inequality to $w^R$.

(4) By definition an exact witness of dimension $d=Q_s(w)$ has
$U(w)e_1=\lambda e_k$ with $|\lambda|=1$, so $\rho(w)=1$, while for
$z\ne w$ of the same length $U(z)e_1\notin\C e_k$. If $\rho(z)=1$ for
such $z$, then $|\langle e_k|U(z)e_1\rangle|=1=\|U(z)e_1\|$, and
equality in Cauchy--Schwarz forces $U(z)e_1\in\C e_k$, a
contradiction; hence $\rho(z)<1=\rho(w)$ for all $z\ne w$ and the gap
is positive.
\end{proof}

\section{Toolbox}\label{sec:toolbox}

This section collects the elementary lemmas used in the main
arguments. Throughout, constants are explicit.

\begin{lemma}[Scalar bounds]\label{lem:scalar}
\emph{(a)} For all real $\beta\ge0$ and integers $n\ge0$,
$(1+\beta)^n\le e^{n\beta}$.
\emph{(b)} For all real $t$,
$|e^{2\pi it}-1|=2|\sin\pi t|\le2\pi\,\dz{t}$.
\end{lemma}

\begin{proof}
(a) $(1+\beta)^n=\sum_{j=0}^n\binom nj\beta^j$ and
$\binom nj\le n^j/j!$, so $(1+\beta)^n\le\sum_{j\ge0}(n\beta)^j/j!
=e^{n\beta}$.
(b) $|e^{2\pi it}-1|^2=(\cos2\pi t-1)^2+\sin^22\pi t
=2-2\cos2\pi t=4\sin^2\pi t$. Both sides are invariant under
$t\mapsto t+1$ and $t\mapsto-t$, so we may assume
$0\le t\le\tfrac12$, where $t=\dz t$ and $\sin\pi t\le\pi t$.
\end{proof}

\begin{lemma}[Born rule is $2$-Lipschitz]\label{lem:born}
Let $x,y\in\C^d$ with $\|x\|,\|y\|\le1$ and let $e$ be a unit
vector. Then
$\bigl|\,|\langle e|x\rangle|^2-|\langle e|y\rangle|^2\,\bigr|
\le2\|x-y\|$.
\end{lemma}

\begin{proof}
Write $a=\langle e|x\rangle$, $b=\langle e|y\rangle$; then
$|a|,|b|\le1$ and
$\bigl||a|^2-|b|^2\bigr|=(|a|+|b|)\bigl||a|-|b|\bigr|
\le2|a-b|=2|\langle e|x-y\rangle|\le2\|x-y\|$
by the triangle and Cauchy--Schwarz inequalities.
\end{proof}

\begin{lemma}[Sphere partition]\label{lem:partition}
Let $D\ge1$ and $0<\varepsilon\le1$. The unit sphere
$S=\{x\in\R^D:\|x\|=1\}$ can be partitioned into
$N\le(5/\varepsilon)^D$ sets, each of diameter at most $\varepsilon$.
\end{lemma}

\begin{proof}
Let $\{x_1,\dots,x_N\}\subseteq S$ be a maximal subset whose points
have pairwise distances $>\varepsilon/2$. The open balls
$B(x_i,\varepsilon/4)$ are pairwise disjoint and contained in the ball
$B(0,1+\varepsilon/4)$; comparing volumes in $\R^D$,
\[
  N\Bigl(\frac\varepsilon4\Bigr)^D\le\Bigl(1+\frac\varepsilon4\Bigr)^D,
  \qquad\text{so}\qquad
  N\le\Bigl(1+\frac4\varepsilon\Bigr)^D\le\Bigl(\frac5\varepsilon\Bigr)^D,
\]
using $\varepsilon\le1$ in the last step. By maximality, every point
of $S$ is within $\varepsilon/2$ of some $x_i$. Define
$C_i=\{x\in S:\ i\text{ is the least index with }
\|x-x_i\|=\min_j\|x-x_j\|\}$; the $C_i$ partition $S$, each
$x\in C_i$ satisfies $\|x-x_i\|\le\varepsilon/2$, and hence
$\operatorname{diam}C_i\le\varepsilon$.
\end{proof}

\begin{lemma}[Simultaneous Dirichlet approximation]\label{lem:dirichlet}
Let $x_1,\dots,x_D\in\R$ and let $Q\ge1$ be an integer. There is an
integer $p$ with $1\le p\le Q^D$ and $\dz{px_j}\le1/Q$ for all
$j=1,\dots,D$.
\end{lemma}

\begin{proof}
For $p=0,1,\dots,Q^D$ consider the fractional-part vectors
$v_p=(\{px_1\},\dots,\{px_D\})\in[0,1)^D$. Partition $[0,1)^D$ into
$Q^D$ boxes $\prod_j[a_j/Q,(a_j+1)/Q)$, $a_j\in\{0,\dots,Q-1\}$. By
pigeonhole two of the $Q^D+1$ vectors, say $v_{p_1}$ and $v_{p_2}$
with $p_1<p_2$, lie in the same box, so
$|\{p_2x_j\}-\{p_1x_j\}|<1/Q$ for every $j$. Put $p=p_2-p_1\in
[1,Q^D]$. For each $j$, $px_j$ differs from
$\{p_2x_j\}-\{p_1x_j\}$ by an integer, whence
$\dz{px_j}\le|\{p_2x_j\}-\{p_1x_j\}|<1/Q$.
\end{proof}

Next, the combinatorics-on-words input.

\begin{lemma}[Lyndon--Sch\"utzenberger; commuting words]\label{lem:LS}
\emph{(a)} If nonempty strings $x,y$ satisfy $xy=yx$, then there are a
string $u$ and integers $s,t\ge1$ with $x=u^s$ and $y=u^t$.
\emph{(b)} Every nonempty string $x$ has a unique primitive string $u$
such that $x=u^s$ for some $s\ge1$ (the \emph{primitive root} of
$x$).
\emph{(c)} If $x,y$ are strings of the same length with $xy=yx$, then
$x=y$.
\end{lemma}

\begin{proof}
See \cite{LS62} for (a) and \cite[Section~1.3]{Lothaire} for both; (b)
follows from (a): if $x=u^s=v^t$ with $u,v$ primitive, then $u$ and
$v$ commute with $x$, hence, by (a) applied twice, are powers of a
common string, and primitivity forces $u=v$. For (c): if $x,y$ are
empty there is nothing to prove; otherwise (a) gives $x=u^s$, $y=u^t$
with $u$ nonempty, and $|x|=|y|$ forces $s=t$.
\end{proof}

Finally, the dictionary between $\SU(2)$ and rotations of the sphere,
and the one fact about rotations that drives Section~\ref{sec:all}.
Let $\sigma_1,\sigma_2,\sigma_3$ be the Pauli matrices and, for
$a\in\R^3$, write $a\cdot\sigma=\sum_ja_j\sigma_j$; every Hermitian
$2\times2$ matrix of trace $1$ has the form $\frac12(I+a\cdot\sigma)$
for a unique $a\in\R^3$, and it is a rank-one projection
$|f\rangle\langle f|$ ($f$ a unit vector) if and only if $\|a\|=1$. The
vector $a=\hat f$ is the \emph{Bloch vector} of $f$; it depends only
on the ray $\C f$.

\begin{lemma}[Rotations and rays]\label{lem:bloch}
\begin{enumerate}
\item For each $U\in\SU(2)$ there is a unique $R_U\in\SO(3)$ with
  $U(a\cdot\sigma)U^\dagger=(R_Ua)\cdot\sigma$ for all $a\in\R^3$;
  the map $U\mapsto R_U$ is a group homomorphism onto $\SO(3)$ with
  kernel $\{\pm I\}$.
\item For a unit vector $f\in\C^2$ and $U\in\SU(2)$: $Uf\in\C f$ if
  and only if $R_U\hat f=\hat f$.
\item A rotation $R\in\SO(3)$ with $R\ne I$ fixes exactly two unit
  vectors, $\pm a$ (its \emph{axis}); and if two rotations $R,R'\ne I$
  have a common fixed unit vector then $RR'=R'R$.
\end{enumerate}
\end{lemma}

\begin{proof}
(1) is the standard double cover $\SU(2)\to\SO(3)$ (see e.g.\
\cite[Ch.~2]{Wagon}): conjugation by $U$ preserves the real vector
space of traceless Hermitian matrices $\{a\cdot\sigma\}$ and the
quadratic form $\det(a\cdot\sigma)=-\|a\|^2$, so it acts by an
orthogonal map $R_U$, which is a rotation by connectedness of
$\SU(2)$; the map is a homomorphism because conjugation is, its image
is all of $\SO(3)$, and $R_U=I$ iff $U$ commutes with all
$a\cdot\sigma$, i.e.\ $U=\pm I$.

(2) $Uf\in\C f$ iff $U|f\rangle\langle f|U^\dagger=|f\rangle\langle
f|$. Writing $|f\rangle\langle f|=\frac12(I+\hat f\cdot\sigma)$, the
left side is $\frac12(I+(R_U\hat f)\cdot\sigma)$ by (1), and two
matrices $\frac12(I+a\cdot\sigma)$, $\frac12(I+a'\cdot\sigma)$ are
equal iff $a=a'$.

(3) A rotation $R\ne I$ of $\R^3$ has $1$ as a simple eigenvalue (its
eigenvalues are $1,e^{\pm i\theta}$ with $\theta\not\equiv0$), so its
fixed unit vectors are $\pm a$ for a unit eigenvector $a$. If $R,R'\ne
I$ both fix $a$, then both restrict to rotations of the plane
$a^\perp$, and rotations of a plane commute; as both fix $a$ as well,
$RR'=R'R$.
\end{proof}

\begin{lemma}[A free pair of rotations]\label{lem:free}
Let $R_0,R_1\in\SO(3)$ be the rotations by the angle $\arccos\frac13$
about the $x$-axis and about the $z$-axis, respectively. Then
$R_0,R_1$ generate a free group of rank two: no nonempty reduced word
in $R_0^{\pm1},R_1^{\pm1}$ equals $I$.
\end{lemma}

\begin{proof}
This is \'Swierczkowski's theorem \cite{Sw58}; a complete elementary
proof is in \cite[Theorem~2.1]{Wagon}. (One shows by induction on the
length of a reduced word that it maps $(0,1,0)^T$ to a vector of the
form $3^{-k}(a\sqrt2,b,c\sqrt2)$ with $a,b,c\in\Z$ and $b\not\equiv0
\pmod3$, so the word is not the identity.)
\end{proof}

\section{All strings have complexity two}\label{sec:all}

For $R_0,R_1\in\SO(3)$ and $z\in\{0,1\}^n$ write, in parallel with the
automaton convention, $R_z=R_{z_n}\cdots R_{z_1}$.

\begin{lemma}\label{lem:comm}
Let $R_0,R_1$ generate a free group of rank two. If $z\ne w$ are
binary strings of the same length $n\ge1$, then $R_z\ne I$, $R_w\ne
I$, and $R_zR_w\ne R_wR_z$.
\end{lemma}

\begin{proof}
A nonempty positive word in $R_0,R_1$ (no inverses) is a nonempty
reduced word, hence $\ne I$ by freeness; this gives $R_z,R_w\ne I$.
Suppose $R_zR_w=R_wR_z$. Both sides are positive words of length
$2n$, spelled by the strings $wz$ and $zw$ respectively (recall that
the first letter of a string is the rightmost factor). In a free group
two reduced words that represent the same element are identical
letter by letter, so $wz=zw$ in the free monoid $\{0,1\}^*$, and
Lemma~\ref{lem:LS}(c) gives $z=w$.
\end{proof}

\begin{theorem}\label{thm:all}
Let $R_0,R_1\in\SO(3)$ generate a free group of rank two (for
instance the pair of Lemma~\ref{lem:free}), and let $U_0,U_1\in\SU(2)$
be any lifts, $R_{U_b}=R_b$. Let $w\in\{0,1\}^n$, $n\ge1$, let $f$ be
a unit eigenvector of $U(w)$, and let $V\in\SU(2)$ satisfy
$Ve_1=f$. Then $(V^\dagger U_0V,\;V^\dagger U_1V,\;k=1)$ is a
$2$-dimensional quantum witness for $w$ with $\rho(w)=1$.
Consequently
\[
  \AQ(w)=2\qquad\text{for every nonempty binary string }w .
\]
\end{theorem}

\begin{proof}
Since $U\mapsto R_U$ is a homomorphism, $R_{U(z)}=R_{z_n}\cdots
R_{z_1}=R_z$ for every string $z$. By Lemma~\ref{lem:comm}, $R_w\ne
I$, so by Lemma~\ref{lem:bloch}(3) $R_w$ has an axis $\pm a$, and
$U(w)$---which has two orthogonal unit eigenvectors, as a unitary
$2\times2$ matrix---has by Lemma~\ref{lem:bloch}(2) exactly the
eigenvectors with Bloch vector $\pm a$. So $\hat f=\pm a$ and
$U(w)f=\lambda f$ with $|\lambda|=1$, whence
$\rho(w)=|\langle f|U(w)f\rangle|^2=1$ for the triple
$(U_0,U_1)$ with initial and accepting vector $f$.

Let $z\ne w$ have length $n$. If $U(z)f\in\C f$ then, by
Lemma~\ref{lem:bloch}(2), $R_z$ fixes $\hat f=\pm a$; as $R_z\ne I$
(Lemma~\ref{lem:comm}), $R_z$ and $R_w$ are rotations $\ne I$ with a
common fixed unit vector, so they commute by
Lemma~\ref{lem:bloch}(3), contradicting Lemma~\ref{lem:comm}. Hence
$U(z)f\notin\C f$, and by the equality case of Cauchy--Schwarz,
$|\langle f|U(z)f\rangle|<1$.

Finally, conjugation by $V$ moves $f$ to $e_1$ without changing any
amplitude: with $U'_b=V^\dagger U_bV$ we have $U'(z)=V^\dagger U(z)V$
and $\langle e_1|U'(z)e_1\rangle=\langle Ve_1|U(z)Ve_1\rangle=\langle
f|U(z)f\rangle$. Thus for the triple $(U'_0,U'_1,1)$ we have
$\rho(w)=1>\rho(z)$ for every $z\ne w$ of length $n$, i.e.\
$\gapp(w)>0$, and $\AQ(w)\le2$. The lower bound is
Proposition~\ref{prop:basic}(1).
\end{proof}

\begin{remark}[Universality and explicitness]\label{rem:universal}
The pair $(U_0,U_1)$ in Theorem~\ref{thm:all} does not depend on $w$:
one fixed pair of $2\times2$ unitaries witnesses every string of every
length, and only the initial state $f$---the axis of $R_w$---changes.
With the rotations of Lemma~\ref{lem:free} the lifts can be taken as
$U_0=\cos\frac\theta2\,I-i\sin\frac\theta2\,\sigma_1$ and
$U_1=\cos\frac\theta2\,I-i\sin\frac\theta2\,\sigma_3$,
$\theta=\arccos\frac13$, so $\cos\frac\theta2=\sqrt{2/3}$ and
$\sin\frac\theta2=\sqrt{1/3}$: all entries lie in $\Q(\sqrt2,\sqrt3,i)$,
and $f$ is an eigenvector of a product of $n$ such matrices. An
exhaustive computation with this pair confirms $\rho(w)=1>\max_{z\ne
w}\rho(z)$ for all $2^n$ strings of every length $n\le12$; the
smallest gap over all strings of length $n$ is approximately
$3.3\cdot10^{-1}$, $3.3\cdot10^{-1}$, $7.3\cdot10^{-2}$,
$4.0\cdot10^{-2}$, $2.7\cdot10^{-4}$, $6.7\cdot10^{-4}$,
$3.0\cdot10^{-5}$, $6.7\cdot10^{-5}$, $3.8\cdot10^{-5}$,
$2.4\cdot10^{-7}$, $1.6\cdot10^{-7}$, $3.0\cdot10^{-8}$ for
$n=1,\dots,12$. This rapid decay is forced: see
Section~\ref{sec:delta}.
\end{remark}

\begin{remark}[Alphabets]\label{rem:alphabets}
Definition~\ref{def:AQ} extends verbatim to an alphabet of any size
$m\ge1$, and so does Theorem~\ref{thm:all}: $\SO(3)$ contains free
subgroups of every finite rank (a free group of rank two contains
free subgroups of all finite ranks, e.g.\ generated by
$R_0^jR_1R_0^{-j}$, $1\le j\le m$), and Lemma~\ref{lem:comm} only uses
that positive words of equal length that commute are equal, which
holds over every alphabet. Thus every nonempty string over every
finite alphabet has maximum-likelihood quantum automatic complexity
exactly two.
\end{remark}

\begin{remark}[Comparison with exact acceptance]\label{rem:lock}
The witness of Theorem~\ref{thm:all} is exact in the sense of
\cite{KH17}---$w$ is accepted with certainty and every other string of
its length with probability $<1$---but its accepting state coincides
with its initial state ($k=1$). In the lock metaphor of \cite{KH17}
the lock starts \emph{unlocked}, is scrambled by every wrong code, and
is restored exactly by the right one. Insisting on $k=2$ (start
locked, end unlocked) changes the geometry: $R_w$ must then be a
half-turn and $\hat f$ must lie on its equator, and one checks in the
same way that such a witness exists if and only if some pair
$(R_0,R_1)$ has $R_w$ a half-turn and $R_z\ne R_w$ for all $z\ne w$
of the same length. That is a codimension-one condition on the pair,
outside the reach of the free-group argument; whether it can always be
met, i.e.\ whether the exact complexities of \cite{KH17} also collapse
to $2$, we leave open.
\end{remark}

\begin{remark}
The contrast with the classical case is stark: $2$-state probabilistic
witnesses exist only for the thin regular family classified in
\cite[Theorem~4.1]{Gill24}, and $\AP$ attains the value $3$, whereas
quantumly, dimension $2$ suffices universally. The geometric reason
is that the pure-state space of a $2$-state PFA is a $1$-dimensional
simplex on which the dynamics is monotone, whereas the pure-state
space of a qubit is a $2$-sphere on which unitaries act as rotations,
and a free group of rotations moves every point to $2^n$ pairwise
distinct places. Hence the complexity information of the measure-once
model resides entirely in the gap structure, to which we now turn.
\end{remark}

\section{Fixed gaps: \texorpdfstring{$A_{Q,\delta}$}{A(Q,delta)} is
unbounded}\label{sec:delta}

The universal witnesses of Theorem~\ref{thm:all} have gaps that
decay rapidly in $n$ (Remark~\ref{rem:universal}). We now show this
is forced, for every witness in any fixed dimension. Recall that the Thue--Morse word is
overlap-free, hence cube-free, and so are all its prefixes
\cite[Ch.~1]{AS03}; thus cube-free binary strings of every length
exist.

Fix a $d$-dimensional witness datum $(U_0,U_1,k)$ and a string
$w\in\{0,1\}^n$, and define the \emph{prefix states}
\[
  \psi_0=e_1,\qquad \psi_m=U_{w_m}\psi_{m-1}
  \quad(1\le m\le n),
\]
so that $\psi_m=U(w_{[1,m]})e_1$ and $\|\psi_m\|=1$ for all $m$.

\begin{lemma}[Block swap]\label{lem:qswap}
Let $0\le i<j<l\le n$ be indices with
\[
  \|\psi_i-\psi_j\|\le\varepsilon,\qquad
  \|\psi_j-\psi_l\|\le\varepsilon,\qquad
  \|\psi_i-\psi_l\|\le\varepsilon .
\]
Write $B_1=w_{(i,j]}$, $B_2=w_{(j,l]}$, $S=w_{(l,n]}$ (so
$w=w_{[1,i]}B_1B_2S$ with $B_1,B_2$ nonempty), and let
$z=w_{[1,i]}B_2B_1S\in\{0,1\}^n$. Then:
\begin{enumerate}
\item $|\rho(z)-\rho(w)|\le6\varepsilon$;
\item $z=w$ if and only if $B_1B_2=B_2B_1$.
\end{enumerate}
\end{lemma}

\begin{proof}
(1) By the recursive definition of $U(\cdot)$ and of the prefix
states, $U(B_1)\psi_i=\psi_j$ and $U(B_2)\psi_j=\psi_l$, and
\[
  U(w)e_1=U(S)\,U(B_2)U(B_1)\,\psi_i,
  \qquad
  U(z)e_1=U(S)\,U(B_1)U(B_2)\,\psi_i .
\]
Since unitaries are isometries,
\[
  U(B_2)\psi_i=U(B_2)\psi_j+U(B_2)(\psi_i-\psi_j)
  =\psi_l+d_1,\qquad
  \|d_1\|=\|\psi_i-\psi_j\|\le\varepsilon,
\]
and therefore
\[
  U(B_1)U(B_2)\psi_i
  =U(B_1)\psi_i+U(B_1)(\psi_l-\psi_i)+U(B_1)d_1
  =\psi_j+U(B_1)(\psi_l-\psi_i)+U(B_1)d_1 .
\]
Comparing with $U(B_2)U(B_1)\psi_i=\psi_l$ and using the triangle
inequality and isometry again,
\[
  \bigl\|U(B_1)U(B_2)\psi_i-U(B_2)U(B_1)\psi_i\bigr\|
  \le\|\psi_j-\psi_l\|+\|\psi_l-\psi_i\|+\|\psi_i-\psi_j\|
  \le3\varepsilon .
\]
Applying the isometry $U(S)$ and Lemma~\ref{lem:born} gives
$|\rho(z)-\rho(w)|\le2\cdot3\varepsilon=6\varepsilon$.

(2) $z$ and $w$ have the common prefix $w_{[1,i]}$ and the common
suffix $S$, and between them the factors $B_2B_1$ and $B_1B_2$, which
have the same length $l-i$; hence $z=w$ iff $B_2B_1=B_1B_2$.
\end{proof}

\begin{theorem}\label{thm:cubefree}
Let $\delta\in(0,1]$ and let $w\in\{0,1\}^n$ be cube-free. Then every
$d$-dimensional quantum witness for $w$ with $\gapp(w)\ge\delta$
satisfies
\[
  n+1\;\le\;3\,\Bigl(\frac{35}{\delta}\Bigr)^{2d},
  \qquad\text{hence}\qquad
  d\;\ge\;\frac{\log\bigl((n+1)/3\bigr)}{2\log(35/\delta)} .
\]
In particular $A_{Q,\delta}$ is unbounded along the Thue--Morse
prefixes, for every fixed $\delta>0$.
\end{theorem}

\begin{proof}
The unit sphere of $\C^d$ is the unit sphere of $\R^{2d}$ under the
standard identification, which is norm-preserving. By
Lemma~\ref{lem:partition} with $D=2d$ and $\varepsilon=\delta/7\le1$,
partition it into $N\le(35/\delta)^{2d}$ cells of diameter at most
$\delta/7$. The $n+1$ prefix states $\psi_0,\dots,\psi_n$ each lie in
some cell. We claim no cell contains $\psi_{i}$ for four distinct
indices; then $n+1\le3N$, which is the assertion.

First, if $i<j<l$ lie in one cell, then the hypotheses of
Lemma~\ref{lem:qswap} hold with $\varepsilon=\delta/7$, so the swapped
string $z$ satisfies $|\rho(z)-\rho(w)|\le6\delta/7<\delta$. If
$z\ne w$ this contradicts $\gapp(w)\ge\delta$; hence $z=w$, i.e.\
$B_1B_2=B_2B_1$ by Lemma~\ref{lem:qswap}(2), and by
Lemma~\ref{lem:LS}(a) the nonempty blocks $B_1,B_2$ are powers of a
common string, which we may take to be the primitive root $u$ of
$B_2$ (Lemma~\ref{lem:LS}(b)).

Now suppose $i_1<i_2<i_3<i_4$ lie in one cell, with consecutive
blocks $B_1=w_{(i_1,i_2]}$, $B_2=w_{(i_2,i_3]}$,
$B_3=w_{(i_3,i_4]}$, all nonempty. Applying the previous paragraph to
the triple $(i_1,i_2,i_3)$: $B_1,B_2\in\{u^s:s\ge1\}$ where $u$ is
the primitive root of $B_2$. Applying it to the triple
$(i_2,i_3,i_4)$: $B_2,B_3\in\{v^s:s\ge1\}$ where $v$ is the primitive
root of $B_2$---the \emph{same} $B_2$---so $v=u$ by uniqueness
(Lemma~\ref{lem:LS}(b)). Hence
$w_{(i_1,i_4]}=B_1B_2B_3=u^{s_1+s_2+s_3}$ with each $s_i\ge1$, so
$uuu$ is a factor of $w$: a cube, contradicting cube-freeness.
\end{proof}

Classically the periodicity escape is real: $A_{P,\delta}(0^n)=2$ for
all $n$ and all $\delta\le1$, via a two-state automaton with an
absorbing reject state \cite{KHpfa}. Quantum mechanically there are no
absorbing states, and recurrence closes the escape:

\begin{theorem}\label{thm:zeros}
Let $\delta\in(0,1]$. Every $d$-dimensional quantum witness for $0^n$
with $\gapp(0^n)\ge\delta$ satisfies
\[
  n\;<\;\Bigl(\frac{52}{\delta}\Bigr)^{2d},
  \qquad\text{hence}\qquad
  A_{Q,\delta}(0^n)\;\ge\;\frac{\log n}{2\log(52/\delta)} .
\]
\end{theorem}

\begin{proof}
Let $e^{i\theta_1},\dots,e^{i\theta_d}$ be the eigenvalues of $U_0$
and $e^{i\theta_{d+1}},\dots,e^{i\theta_{2d}}$ those of $U_1$. Apply
Lemma~\ref{lem:dirichlet} to the $2d$ numbers $x_j=\theta_j/2\pi$ with
$Q=\lceil16\pi/\delta\rceil$: there is an integer $p$ with
$1\le p\le Q^{2d}$ and $\dz{p\theta_j/2\pi}\le1/Q$ for all $j$. By the
diagonalization identity of Section~\ref{sec:notation} and
Lemma~\ref{lem:scalar}(b),
\[
  \|U_b^p-I\|=\max_j\bigl|e^{ip\theta_j}-1\bigr|
  \le2\pi\max_j\dz{p\theta_j/2\pi}\le\frac{2\pi}{Q}
  \qquad(b=0,1).
\]
Suppose, for contradiction, that $n\ge p$. The string
$z=0^{n-p}1^p$ differs from $0^n$ (as $p\ge1$), and by the reading
convention $U(0^n)e_1=U_0^{\,p}\psi$ and $U(z)e_1=U_1^{\,p}\psi$,
where $\psi=U_0^{\,n-p}e_1$ is a unit vector. By
Lemma~\ref{lem:born} and the triangle inequality,
\[
  \bigl|\rho(0^n)-\rho(z)\bigr|
  \le2\,\bigl\|(U_0^{\,p}-U_1^{\,p})\psi\bigr\|
  \le2\bigl(\|U_0^{\,p}-I\|+\|I-U_1^{\,p}\|\bigr)
  \le\frac{8\pi}{Q}\le\frac{8\pi\,\delta}{16\pi}=\frac\delta2<\delta,
\]
contradicting $\gapp(0^n)\ge\delta$. Hence $n<p\le Q^{2d}$, and
$Q\le16\pi/\delta+1\le(16\pi+1)/\delta\le52/\delta$ using
$\delta\le1$.
\end{proof}

\begin{corollary}\label{cor:unbounded}
For every fixed $\delta>0$, $A_{Q,\delta}$ is unbounded, both along
the cube-free strings and along the strings $0^n$.
\end{corollary}

\begin{remark}[Quantum locks wear out]\label{rem:wear}
Theorem~\ref{thm:zeros} is an approximate-analysis sibling of the
Jordan--Schur mechanism used in \cite{KH17} to prove the
finite-group complexity $Q_f$ unbounded: there, powers of elements of
a finite unitary group eventually commute exactly; here, powers of
\emph{arbitrary} unitaries recur near the identity at a rate
controlled only by the dimension. In the lock metaphor of
\cite{KH17}: no fixed-dimensional quantum lock can maintain a fixed
security margin against arbitrarily long codes---not even against the
all-zeros code. The hypothesis of Theorem~\ref{thm:zeros} is tight in
the sense that the margin, not the lock, must decay:
$A_{Q,1}(0^n)\le n+1$, via $d=n+1$, $U_0$ the cyclic permutation
$e_i\mapsto e_{i+1\bmod(n+1)}$, $U_1=I$, and accepting state
$e_{n+1}$; then $U(z)e_1=e_{1+j\bmod(n+1)}$ where $j$ is the number of
zeros in $z$, so $\rho(z)=1$ if $z=0^n$ and $\rho(z)=0$ for every
other $z$ of length $n$, giving gap $1$. Numerically, the largest gap
achievable for $0^n$ in dimension $2$ decays steadily
($\approx0.53,\,0.36,\,0.28,\,0.18,\,0.16$ at
$n=4,6,10,12,14$), in accordance with the theorem.
\end{remark}

\begin{remark}\label{rem:allwords}
The two mechanisms plausibly combine. On one side, the proof of
Theorem~\ref{thm:cubefree} shows more generally that if $w$ contains
no $E$-th power ($E\ge3$), then no cell contains more than $E$ prefix
states, giving $n+1\le E(35/\delta)^{2d}$ by the same argument with
$E+1$ indices in place of four. On the other side, long powers invite
the recurrence argument applied to $U(u)$ for the repeated block
$u$. We expect, but have not proved, that
$A_{Q,\delta}(w)=\Omega_\delta(\log|w|)$ for \emph{every} string $w$;
the missing case is a string containing powers of intermediate
exponent. If so, then in contrast with $A_{P,\delta}$---which
separates periodic from aperiodic strings by staying bounded on the
former---the quantum measure grows on everything, and the
discriminating quantity becomes the \emph{decay rate} of the best gap
at fixed dimension. Our experiments are consistent with a stark
difference in rate: gently polynomial for $0^n$, much faster for
aperiodic strings.
\end{remark}

\begin{question}\label{q:rate}
What is the growth rate of $A_{Q,\delta}$ along the Thue--Morse
prefixes, and along $0^n$, for fixed $\delta$? Both are between
$\Omega_\delta(\log n)$ and $O(n^2)$
(Proposition~\ref{prop:quadupper}; for $0^n$, indeed $O(n)$ by
Remark~\ref{rem:wear}), and we do not know whether a fixed gap can be
maintained in polylogarithmic dimension on any infinite family of
strings.
\end{question}

\begin{remark}
Like Theorem~\ref{thm:all} (Remark~\ref{rem:alphabets}), the results
of this section extend verbatim to arbitrary finite alphabets; over a
ternary alphabet, square-free strings give Theorem~\ref{thm:cubefree}
with $2$ points per cell in place of $3$, using square-freeness in
place of cube-freeness.
\end{remark}

\section{Computational aspects}\label{sec:complexity}

The preceding sections concern the \emph{values} of $\AQ$ and
$A_{Q,\delta}$; we now locate their computational difficulty. In
outline: the string-maximization problem underlying witness
verification is NP-complete, with inverse-polynomial robustness and
exact rational acceptance values (Theorem~\ref{thm:nphard});
verification of $\delta$-witnesses is coNP-complete
(Corollary~\ref{cor:conp}); the synthesis promise problem lies in
$\Sigma_2^p$ (Proposition~\ref{prop:sigma2}); and the exact value
$A_{Q,\delta}(w)$ is computable in exponential time, via quantifier
elimination over the reals together with a quadratic universal upper
bound (Propositions~\ref{prop:quadupper} and~\ref{prop:exact}).

Throughout this section, an \emph{exact instance} presents the entries
of $U_0,U_1$ as elements of a number field of degree $O(1)$ over
$\Q$, by rational coordinates of polynomial bit-size in a fixed
basis; a \emph{rational instance} presents rational entries. The
dimension $d$ and the length $n$ are given in unary.

\begin{lemma}[Clocking]\label{lem:clock}
Let $W_{b,m}\in U(B)$ be given for $b\in\{0,1\}$ and $1\le m\le n$.
On $\C^{n+1}\otimes\C^B$ define
\[
  U_b \;=\; \sum_{m=1}^{n} |m{+}1\rangle\langle m| \otimes W_{b,m}
  \;+\; |1\rangle\langle n{+}1| \otimes I .
\]
Then $U_0,U_1$ are unitary, and with start state $|1\rangle\otimes
e_1$ and accepting state $|n{+}1\rangle\otimes e_{k'}$, every
$z\in\{0,1\}^n$ satisfies
\[
  \rho(z)=\bigl|\langle e_{k'}|\,W_{z_n,n}\cdots W_{z_1,1}\,
  e_1\rangle\bigr|^2 .
\]
Thus an MO-QFA can execute an arbitrary position-dependent gate
sequence, at a factor $n+1$ in dimension.
\end{lemma}

\begin{proof}
Unitarity: $U_b^\dagger U_b=\sum_{m=1}^n|m\rangle\langle m|\otimes
W_{b,m}^\dagger W_{b,m}+|n{+}1\rangle\langle n{+}1|\otimes I=I$, since
the clock summands act on orthogonal blocks. By induction on $m$, the
state after reading $z_1\cdots z_m$ is $|m{+}1\rangle\otimes
W_{z_m,m}\cdots W_{z_1,1}e_1$: the wrap-around term
$|1\rangle\langle n{+}1|$ is never activated on inputs of length
exactly $n$.
\end{proof}

\begin{proposition}[Quadratic upper bound]\label{prop:quadupper}
For every $w\in\{0,1\}^n$ and every $\delta\in(0,1]$,
\[
  A_{Q,\delta}(w)\;\le\;A_{Q,1}(w)\;\le\;(n+1)^2 .
\]
\end{proposition}

\begin{proof}
Apply Lemma~\ref{lem:clock} with work register $\C^{n+1}$, basis
$|0\rangle,\dots,|n\rangle$, cyclic increment
$C|c\rangle=|c+1\bmod (n{+}1)\rangle$, and
$W_{b,m}=C^{[\,b\ne w_m\,]}$: a mismatch counter. After reading $z$
the work state is $|d_H(z,w)\bmod(n{+}1)\rangle=|d_H(z,w)\rangle$,
where $d_H$ is Hamming distance and $d_H\le n$. With accepting state
$|n{+}1\rangle\otimes|0\rangle$ we get $\rho(z)=1$ if $z=w$ and
$\rho(z)=0$ otherwise: gap $1$ in dimension $(n+1)^2$. (For $0^n$,
Remark~\ref{rem:wear} improves this to $n+1$.)
\end{proof}

\begin{lemma}[Stability]\label{lem:stability}
Let $U_0,U_1\in U(d)$, let $\widetilde U_0,\widetilde U_1$ be
arbitrary matrices with $\|\widetilde U_b-U_b\|\le\tau\le\frac1{2n}$,
and let $\widetilde\rho$ be defined from the $\widetilde U_b$ by the
same formula as $\rho$. Then:
\begin{enumerate}
\item $|\widetilde\rho(z)-\rho(z)|\le6n\tau$ for every
  $z\in\{0,1\}^n$;
\item if moreover $\|\widetilde U_b^\dagger\widetilde U_b-I\|
  \le\varepsilon<1$, then there are unitaries $W_b$ with
  $\|\widetilde U_b-W_b\|\le\varepsilon$.
\end{enumerate}
\end{lemma}

\begin{proof}
(1) Telescoping,
$\widetilde U(z)-U(z)=\sum_{m=1}^n
\widetilde U_{z_n}\cdots\widetilde U_{z_{m+1}}
(\widetilde U_{z_m}-U_{z_m})U_{z_{m-1}}\cdots U_{z_1}$,
and $\|\widetilde U_b\|\le1+\tau$, so each summand has norm at most
$(1+\tau)^{n}\tau\le e^{n\tau}\tau\le e^{1/2}\tau\le2\tau$
(Lemma~\ref{lem:scalar}(a)); hence
$\|\widetilde U(z)e_1-U(z)e_1\|\le2n\tau$. With
$a=\langle e_k|U(z)e_1\rangle$ and
$\tilde a=\langle e_k|\widetilde U(z)e_1\rangle$ we have $|a|\le1$,
$|\tilde a|\le(1+\tau)^n\le2$, and
$\bigl||\tilde a|^2-|a|^2\bigr|\le(|\tilde a|+|a|)|\tilde a-a|
\le3\cdot2n\tau$.
(2) Let $\widetilde U_b=W_bP_b$ be a polar decomposition with $W_b$
unitary and $P_b=(\widetilde U_b^\dagger\widetilde U_b)^{1/2}$
positive (invertible since $\varepsilon<1$). Then
$\|\widetilde U_b-W_b\|=\|P_b-I\|
=\max_\lambda|\sqrt\lambda-1|\le\max_\lambda|\lambda-1|
=\|\widetilde U_b^\dagger\widetilde U_b-I\|\le\varepsilon$,
using $|\sqrt x-1|\le|x-1|$ for $x\ge0$.
\end{proof}

\begin{lemma}[Prescribed first row]\label{lem:householder}
For every real unit vector $r\in\R^D$ there is a real symmetric
orthogonal matrix $H$ (a \emph{Householder reflection}) with
$He_1=r$, hence with first row $r^{T}$; its entries are rational
functions of the entries of $r$.
\end{lemma}

\begin{proof}
If $r=e_1$, take $H=I$. Otherwise set $u=e_1-r\ne0$ and
$H=I-\dfrac{2uu^{T}}{u^{T}u}$. Then $H$ is symmetric, and
\[
  H^2=I-\frac{4uu^{T}}{u^{T}u}+\frac{4u(u^{T}u)u^{T}}{(u^{T}u)^2}=I,
\]
so $H$ is orthogonal. Since $\|r\|=\|e_1\|=1$ we have
$u^{T}u=2-2r_1=2\,u^{T}e_1$, whence
\[
  He_1=e_1-\frac{2u\,(u^{T}e_1)}{u^{T}u}=e_1-u=r ,
\]
and by symmetry the first row of $H$ is $(He_1)^{T}=r^{T}$.
\end{proof}

\subsection*{Maximization is NP-complete}

\begin{definition}
\textup{\textsc{QFA-MaxString}}: given $n$, $d$ (unary), an exact instance
$(U_0,U_1)$, an accepting index $k$, and $\theta\in\Q$: is there
$z\in\{0,1\}^n$ with $\rho(z)\ge\theta$?
\end{definition}

\begin{theorem}\label{thm:nphard}
\textup{\textsc{QFA-MaxString}} is NP-hard, via a polynomial-time reduction
from \textup{\textsc{3-Sat}} \cite{GJ79} under which, for a formula with $J$
clauses,
\[
  \max_{z}\rho(z)=\frac{\mathrm{maxsat}(\varphi)^2}{7J^2}
  \;\in\;\Q ,
\]
so that satisfiable instances attain $\max_z\rho=\frac17$ while
unsatisfiable ones have
$\max_z\rho\le\frac17\bigl(1-\frac1J\bigr)^2
\le\frac17-\frac1{7J}$. Consequently NP-hardness holds even as a
promise problem with inverse-polynomial gap; and the promise variant
with rational instances exponentially close to unitary is
NP-complete.
\end{theorem}

\begin{proof}
Let $\varphi$ have variables $x_1,\dots,x_n$ and clauses
$C_1,\dots,C_J$, each on three distinct variables. Identify strings
$z\in\{0,1\}^n$ with truth assignments; note every string of length
$n$ is a legal assignment. The work register is $\C^{8J}$ with basis
$|j,y\rangle$, $j\in[J]$, $y\in\{0,1\}^3$; block $j$ will store the
values of the three variables of $C_j$.

\emph{Fan-out.} Let $F$ be the Householder reflection of
Lemma~\ref{lem:householder} with
$F|1,000\rangle=\frac1{\sqrt J}\sum_j|j,000\rangle$, taking the
basis ordered with $|1,000\rangle$ first.

\emph{Storage.} For $b\in\{0,1\}$ and $1\le m\le n$ let $P_{b,m}$ act
on block $j$ by flipping stored bit $r$ whenever $b=1$ and $x_m$ is
the $r$-th variable of $C_j$, and trivially otherwise; each $P_{b,m}$
is a permutation matrix, and $P_{0,m}=I$. After reading $z$, block
$j$ holds $|y_j\rangle$ with $y_j=z\!\restriction_{C_j}$, since each
stored bit is toggled exactly when its variable is read with value
$1$. No absorption occurs anywhere---by
Theorem~\ref{thm:zeros} none is available---so the clause
disjunctions are not computed during reading at all; they are scored
by interference at the end.

\emph{Fan-in.} A $3$-clause has exactly $7$ satisfying local
assignments, so
$v_j=\frac1{\sqrt7}\sum_{y\models C_j}|j,y\rangle$ is a unit vector,
and $\frac1{\sqrt J}(v_1^\dagger,\dots,v_J^\dagger)$ is a unit row;
let $G$ be the Householder reflection with this first row
(Lemma~\ref{lem:householder} again).

Apply Lemma~\ref{lem:clock} with $W_{b,1}=P_{b,1}F$,
$W_{b,m}=P_{b,m}$ for $1<m<n$, $W_{b,n}=G\,P_{b,n}$, and accepting
work vector $|1,000\rangle$. The final amplitude is
\[
  \frac1{\sqrt J}\sum_{j=1}^J\frac1{\sqrt J}\langle v_j|y_j\rangle
  =\frac1J\sum_{j=1}^J\frac{[\,z\models C_j\,]}{\sqrt7},
  \qquad\text{whence}\qquad
  \rho(z)=\frac{\mathrm{sat}_\varphi(z)^2}{7J^2}
\]
exactly, for every $z\in\{0,1\}^n$. Taking $\theta=\frac17$ proves
hardness, with the stated yes/no separation
$\frac17-\frac17(1-\frac1J)^2=\frac{2J-1}{7J^2}\ge\frac1{7J}$. The
dimension is $8J(n+1)$, and all entries lie in $\Q(\sqrt7,\sqrt J)$
with polynomial bit-size, so the map is polynomial-time and the
instance is exact; membership in NP for exact instances is immediate
(guess $z$ and evaluate $\rho(z)$ by exact arithmetic in the given
field). For the rational promise variant, replace $F$ and $G$ by
entrywise rational approximations to accuracy $2^{-p}$,
$p=\mathrm{poly}$: by Lemma~\ref{lem:stability} every $\rho(z)$
moves by at most $6n\cdot O(2^{-p}d)$, preserving the $\frac1{7J}$
promise gap, while verification becomes rational arithmetic with the
tolerance absorbed by the promise.
\end{proof}

\subsection*{Witness verification is coNP-complete}

We now show that verifying a claimed $\delta$-witness is
coNP-complete. The reduction refines the machine of
Theorem~\ref{thm:nphard} by a \emph{reference branch}: a reserved
direction of the work register gives one designated string $w$ a
tunable head start, so that the \emph{size of the gap} at
$w$---rather than the maximum of $\rho$---encodes unsatisfiability.

\begin{definition}[The referenced machine]\label{def:refmachine}
Let $\varphi$ be a \textup{\textsc{3-Sat}} formula with $J$ clauses
$C_1,\dots,C_J$, each on three distinct variables, in which every
variable occurs in some clause (variables occurring in no clause may
be deleted without affecting satisfiability). Let
$x_1,\dots,x_{n-1}$ be its variables, let $x_n$ be a fresh variable
occurring in no clause, and let $\mu\in(0,1)$ be a rational
parameter. The machine $M_\varphi(\mu)$ is the clocked machine of
Lemma~\ref{lem:clock} on the work register $\C^{8J+1}$, with basis
$|j,y\rangle$ ($j\in[J]$, $y\in\{0,1\}^3$) together with one
reserved vector $|\mathrm{ref}\rangle$, built from the following
gates. All Householder reflections below exist by
Lemma~\ref{lem:householder}, taking the basis ordered with the
relevant vector first.
\begin{itemize}
\item \emph{Fan-out.} $F'$ is the Householder reflection with
\[
  F'\,|1,000\rangle
  =\sqrt{1-\mu}\;\frac1{\sqrt J}\sum_{j=1}^J|j,000\rangle
  +\sqrt{\mu}\;|\mathrm{ref}\rangle ,
\]
a legitimate target since its squared norm is $(1-\mu)+\mu=1$.
\item \emph{Storage.} $P'_{b,m}=P_{b,m}\oplus[1]$, where $P_{b,m}$
is the storage permutation of Theorem~\ref{thm:nphard}; thus the
stored bits are updated exactly as before while
$|\mathrm{ref}\rangle$ is fixed. Since $x_n$ occurs in no clause,
$P'_{b,n}=I$ for both letters.
\item \emph{Two fan-ins.} $G_0$ is the Householder reflection whose
first row is $r_0^{T}$, where
$r_0=\frac1{\sqrt J}\sum_{j}v_j$ with
$v_j=\frac1{\sqrt7}\sum_{y\models C_j}|j,y\rangle$ as in
Theorem~\ref{thm:nphard} (so $r_0$ vanishes on
$|\mathrm{ref}\rangle$); and $G_1$ is the Householder reflection
whose first row is $r_1^{T}$, where
\[
  r_1=\sqrt{\mu}\;\frac1{\sqrt J}\sum_{j=1}^J|j,000\rangle
  +\sqrt{1-\mu}\;|\mathrm{ref}\rangle ,
\]
again a unit vector since $\mu+(1-\mu)=1$.
\end{itemize}
The letter gates are $W_{b,1}=P'_{b,1}F'$, $W_{b,m}=P'_{b,m}$ for
$1<m<n$, $W_{0,n}=G_0P'_{0,n}$ and $W_{1,n}=G_1P'_{1,n}$; the
accepting work vector is $|1,000\rangle$.
\end{definition}

\begin{lemma}[Acceptance function of the referenced
machine]\label{lem:refaccept}
Write each input as $z=ab$ with $a\in\{0,1\}^{n-1}$ and
$b\in\{0,1\}$, and let
\[
  Z(a)\;=\;\#\{\,j\in[J]:\ a\text{ assigns }0\text{ to every
  variable of }C_j\,\}.
\]
Then, for every $a$,
\[
  \rho(a0)=(1-\mu)\,\frac{\mathrm{sat}_\varphi(a)^2}{7J^2},
  \qquad
  \rho(a1)=\mu(1-\mu)\Bigl(1+\frac{Z(a)}J\Bigr)^{2}.
\]
Moreover $Z(a)=J$ if and only if $a=0^{n-1}$.
\end{lemma}

\begin{proof}
By Lemma~\ref{lem:clock} it suffices to track the work register
under $W_{z_n,n}\cdots W_{z_1,1}$. After the fan-out, the state is
$\sqrt{1-\mu}\,\frac1{\sqrt J}\sum_j|j,000\rangle
+\sqrt\mu\,|\mathrm{ref}\rangle$. The storage gates act blockwise
and fix $|\mathrm{ref}\rangle$, and---exactly as in
Theorem~\ref{thm:nphard}---each stored bit is toggled precisely when
its variable is read with value $1$; since $x_n$ occurs in no
clause, the stored contents depend only on $a$. Hence the state
reaching the position-$n$ fan-in is
\[
  \Psi(a)=\sqrt{1-\mu}\;\frac1{\sqrt J}\sum_{j=1}^J|j,y_j\rangle
  +\sqrt{\mu}\;|\mathrm{ref}\rangle ,
  \qquad y_j=a\!\restriction_{C_j}.
\]
If $b=0$, the accepting amplitude is
\[
  \langle r_0|\Psi(a)\rangle
  =\sqrt{1-\mu}\;\frac1{J}\sum_{j=1}^J\langle v_j|j,y_j\rangle
  =\sqrt{1-\mu}\;\frac{\mathrm{sat}_\varphi(a)}{J\sqrt7},
\]
since $\langle v_j|j,y_j\rangle=[\,y_j\models C_j\,]/\sqrt7$ and
$y_j\models C_j$ iff $a\models C_j$, while $r_0$ has no
$|\mathrm{ref}\rangle$ component. If $b=1$, the accepting amplitude
is
\[
  \langle r_1|\Psi(a)\rangle
  =\sqrt{\mu}\,\sqrt{1-\mu}\;\frac{Z(a)}J
  +\sqrt{1-\mu}\,\sqrt{\mu}
  =\sqrt{\mu(1-\mu)}\Bigl(1+\frac{Z(a)}J\Bigr),
\]
since $\langle j,000|j,y_j\rangle=[\,y_j=000\,]$. Squaring gives the
two formulas.

Finally, if $a_i=1$ for some $i$, then $x_i$ occurs in some clause
$C_j$, so block $j$ holds $y_j\ne000$ and $Z(a)\le J-1$; conversely
$a=0^{n-1}$ gives $y_j=000$ for every $j$, so $Z(a)=J$.
\end{proof}

\begin{lemma}[Separation]\label{lem:refsep}
In the referenced machine, take
\[
  \mu:=\tfrac1{28}+\tfrac1{224J}
  \qquad\text{(so that } 4\mu=\tfrac17+\tfrac1{56J}\text{)},
  \qquad w:=0^{n-1}1 .
\]
Then:
\begin{enumerate}
\item $\rho(w)=4\mu(1-\mu)
  =(1-\mu)\bigl(\tfrac17+\tfrac1{56J}\bigr)$;
\item every $a$ satisfies
  $\rho(a0)\le\frac{1-\mu}7=\rho(w)-\frac{1-\mu}{56J}$, with
  equality if and only if $a$ satisfies $\varphi$;
\item if $\varphi$ is unsatisfiable, every $a$ satisfies
  $\rho(a0)\le\rho(w)-\frac{9(1-\mu)}{56J}$;
\item every $a\ne0^{n-1}$ satisfies
  $\rho(a1)\le\rho(w)-\frac{6(1-\mu)}{56J}$.
\end{enumerate}
In particular $w$ is the strict maximizer of $\rho$, for every
$\varphi$.
\end{lemma}

\begin{proof}
(1) By Lemma~\ref{lem:refaccept} with $Z(0^{n-1})=J$,
$\rho(w)=\mu(1-\mu)\cdot2^2$, and
$4\mu(1-\mu)=(1-\mu)\cdot4\mu
=(1-\mu)\bigl(\tfrac17+\tfrac1{56J}\bigr)$.

(2) From $\mathrm{sat}_\varphi(a)\le J$,
$\rho(a0)\le(1-\mu)\frac{J^2}{7J^2}=\frac{1-\mu}7$, with equality
iff $\mathrm{sat}_\varphi(a)=J$; and by (1),
$\rho(w)-\frac{1-\mu}7
=(1-\mu)\bigl(\tfrac17+\tfrac1{56J}-\tfrac17\bigr)
=\frac{1-\mu}{56J}$.

(3) Now $\mathrm{sat}_\varphi(a)\le J-1$, so
$\rho(a0)\le\frac{1-\mu}7\bigl(1-\frac1J\bigr)^2$, and
\[
  \rho(w)-\frac{1-\mu}7\Bigl(1-\frac1J\Bigr)^{2}
  =(1-\mu)\biggl[\frac1{56J}
  +\frac17\Bigl(\frac2J-\frac1{J^2}\Bigr)\biggr]
  \;\ge\;(1-\mu)\Bigl[\frac1{56J}+\frac1{7J}\Bigr]
  =\frac{9(1-\mu)}{56J},
\]
using $\frac2J-\frac1{J^2}\ge\frac1J$, i.e.\
$\frac1J\ge\frac1{J^2}$, valid for $J\ge1$.

(4) By Lemma~\ref{lem:refaccept}, $Z(a)\le J-1$, and $\rho(a1)$ is
increasing in $Z(a)$, so
$\rho(a1)\le\mu(1-\mu)\bigl(2-\frac1J\bigr)^2$. Then
\[
  \rho(w)-\mu(1-\mu)\Bigl(2-\frac1J\Bigr)^{2}
  =\mu(1-\mu)\Bigl(\frac4J-\frac1{J^2}\Bigr)
  \;\ge\;\frac{3\mu(1-\mu)}J
  =\frac{3\,\rho(w)}{4J}
  \;\ge\;\frac3{4J}\cdot\frac{1-\mu}7
  =\frac{6(1-\mu)}{56J},
\]
using first $\frac4J-\frac1{J^2}\ge\frac3J$ ($J\ge1$) and then
$\rho(w)\ge\frac{1-\mu}7$ from (1).

Strict maximality: every string other than $w$ is either of the form
$a0$, covered by (2) with deficit at least $\frac{1-\mu}{56J}>0$, or
of the form $a1$ with $a\ne0^{n-1}$, covered by (4).
\end{proof}

\begin{corollary}\label{cor:conp}
The problem \textup{\textsc{GapVerify}}---given an exact (or
promise-rational) witness datum $(U_0,U_1,k)$, a string $w$, and
$\delta\in\Q$: is $\gapp(w)\ge\delta$?---is coNP-complete.
\end{corollary}

\begin{proof}
Membership: a certificate for the complement is a single $z\ne w$
with $\rho(w)-\rho(z)<\delta$, checkable by exact (respectively
rational) arithmetic.

Hardness, by reduction from the complement of
\textup{\textsc{3-Sat}}: map $\varphi$ (preprocessed as in
Definition~\ref{def:refmachine}) to the instance consisting of
$M_\varphi(\mu)$, the string $w=0^{n-1}1$, and
\[
  \delta:=\frac{5(1-\mu)}{56J},
  \qquad \mu=\tfrac1{28}+\tfrac1{224J} .
\]
If $\varphi$ is satisfiable, a satisfying assignment $a$ attains
equality in Lemma~\ref{lem:refsep}(2), so
$\gapp(w)\le\frac{1-\mu}{56J}<\delta$. If $\varphi$ is
unsatisfiable, then by Lemma~\ref{lem:refsep}(3) and (4) every
$z\ne w$ satisfies
$\rho(z)\le\rho(w)-\frac{6(1-\mu)}{56J}\le\rho(w)-\delta$, so
$\gapp(w)\ge\delta$. By the last clause of
Lemma~\ref{lem:refsep}, the constructed datum is a genuine witness
for $w$ in both cases: only its $\delta$-quality encodes
satisfiability, so the reduction cannot be defeated by rejecting
malformed witnesses.

The map is polynomial-time: $\mu$ is rational, the dimension is
$(8J+1)(n+1)$, and by Lemma~\ref{lem:householder} all entries are
rational functions of $\sqrt\mu$, $\sqrt{1-\mu}$, $\sqrt7$ and
$\sqrt J$, a number field of degree $O(1)$ presented by
polynomial-size rational coordinates. The margins on the two sides
of $\delta$ are at least $\frac{1-\mu}{56J}=\Theta(1/J)$, so the
rationalization of Lemma~\ref{lem:stability} applies verbatim. (The
construction and these parameter choices were verified by machine
on satisfiable and unsatisfiable instances.)
\end{proof}

\subsection*{Upper bounds for synthesis}

\begin{proposition}\label{prop:sigma2}
The promise problem---given $w\in\{0,1\}^n$, $d$ in unary, and
rationals $\delta,\eta$ with $0<\eta\le\delta\le1$: accept if some
$d$-dimensional witness for $w$ has $\gapp\ge\delta$, reject if every
$d$-dimensional unitary triple has $\gapp\le\delta-\eta$---is in
$\Sigma_2^p$. For each fixed $d$ it is in
$\mathrm{P}^{\mathrm{NP}}$ when $\eta\ge1/\mathrm{poly}(n)$.
\end{proposition}

\begin{proof}
Set $\tau=\eta/(60n)$ and let $\Gamma$ be the grid of matrix pairs
whose entries are Gaussian rationals with real and imaginary parts
multiples of $\gamma=\tau/(2d)$ and of modulus at most $2$: elements
of $\Gamma$ have polynomial bit-size. The $\Sigma_2^p$ certificate is
a pair $(\widetilde U_0,\widetilde U_1)\in\Gamma$ together with $k$,
checked as follows in polynomial time with one universal quantifier:
(i) $\|\widetilde U_b^\dagger\widetilde U_b-I\|_F\le2\tau$; (ii) for
all $z\ne w$: $\widetilde\rho(w)-\widetilde\rho(z)\ge\delta-\eta/2$
(a coNP predicate, by exact rational arithmetic). If some true
witness has gap $\ge\delta$, rounding its entries to $\Gamma$
perturbs each $U_b$ by at most $d\gamma=\tau/2$ in Frobenius, hence
operator, norm; then (i) holds and, by
Lemma~\ref{lem:stability}(1), each $\widetilde\rho(z)$ is within
$3n\tau=\eta/20$ of $\rho(z)$, so (ii) holds. Conversely if some
certificate passes, Lemma~\ref{lem:stability}(2) yields unitaries
$W_b$ within $2\tau$ of $\widetilde U_b$, whose gap is at least
$(\delta-\eta/2)-2\cdot12n\tau=\delta-\eta/2-2\eta/5>\delta-\eta$,
so the instance is not a reject instance. For fixed $d$ and
$\eta\ge1/\mathrm{poly}$, the grid has
$(O(1/\gamma))^{4d^2}=\mathrm{poly}(n)$ elements; enumerate it and
decide predicate (ii) for each element with one NP-oracle query.
\end{proof}

\begin{proposition}\label{prop:exact}
For rational $\delta\in(0,1]$, the exact value $A_{Q,\delta}(w)$ is
computable, in deterministic time
$2^{O(n\cdot A_{Q,\delta}(w)^2)}\subseteq2^{O(n^5)}$; likewise
$\AQ(w)$.
\end{proposition}

\begin{proof}
First, the best gap is monotone in $d$: padding $U_b\mapsto
U_b\oplus I_{d'-d}$ preserves every amplitude
$\langle e_k|U(z)e_1\rangle$, so a $d$-dimensional witness yields a
$d'$-dimensional one. Second, for each fixed $d$ the predicate
``some $d$-dimensional witness for $w$ has $\gapp\ge\delta$'' is a
sentence of the existential theory of the reals: $8d^2$ real
variables (real and imaginary parts), $4d^2$ polynomial equations of
degree $2$ expressing $U_b^\dagger U_b=I$, and $2^n-1$ polynomial
inequalities of degree $2n$ expressing the gaps, with coefficients of
bit-size $\mathrm{poly}(n)$. By Renegar's decision procedure
\cite{Ren92} this sentence is decidable in time
$\bigl(2^n\cdot n\bigr)^{O(d^2)}=2^{O(nd^2)}$. Search
$d=2,3,\dots$; by Proposition~\ref{prop:quadupper} the search
terminates by $d=(n+1)^2$, and it stops at $d=A_{Q,\delta}(w)$, for
the stated bound. For $\AQ$, replace the gap inequalities by strict
ones and $\delta$ by $0$; termination via
$\AQ(w)\le A_{Q,1}(w)\le(n+1)^2$.
\end{proof}

\begin{remark}\label{rem:lit}
The classical analogue of Theorem~\ref{thm:nphard}---the most
probable string for a probabilistic automaton---is NP-hard by
Casacuberta and de la Higuera \cite{CH00}. The quantum statement has
some pleasant additional features: every string of length $n$ is a
legal input (no consensus-string padding), the acceptance values are
exact rationals $\mathrm{sat}^2/(7J^2)$ (hardness owes nothing to
numerical precision), and the promise gap is inverse-linear. The
mechanism is the obstruction of Theorem~\ref{thm:zeros} turned from
bug into feature: unable to compute clause disjunctions by
absorption, the machine stores literals reversibly and scores all
clauses simultaneously, by interference, at measurement time. For
unbounded input length, Blondel, Jeandel, Koiran and Portier
\cite{BJKP05} showed that strict-cutpoint emptiness for measure-once
QFAs is decidable while the non-strict version is undecidable;
fixing the length trivializes decidability, and this section locates
the residual complexity.
\end{remark}

\begin{question}\label{q:complexity}
Two questions remain open. (a) Is \textup{\textsc{QFA-MaxString}} hard for
\emph{fixed} dimension---in particular $d=2$, where the acceptance
function is $\rho(z)=|\langle e_k|U(z)e_1\rangle|^2$ for a product of
$2\times2$ matrices and, by Theorem~\ref{thm:all}, every string is
the maximizer for a suitable initial state? The reduction of
Theorem~\ref{thm:nphard} needs $d=8J(n+1)$. (b) What is the
complexity of computing $A_{Q,\delta}(w)$ from the bare input
$(w,\delta)$? No hardness is known; encoding computation into a
\emph{string}, rather than into an automaton, is precisely the
obstacle that has left the complexity of Shallit--Wang automatic
complexity \cite{SW01} open. (By Theorem~\ref{thm:all} the
unparameterized value $\AQ(w)=2$ is trivial to compute.)
\end{question}

\section*{Acknowledgments}

The results in this note were found in conversation with, and verified
computationally by, Claude Fable 5 (Anthropic); the exhaustive check
of Remark~\ref{rem:universal}, and the reduction of
Section~\ref{sec:complexity} together with the parameter choices in
Corollary~\ref{cor:conp}, were machine-verified. The author thanks Kenneth Gill, whose paper \cite{Gill24}
prompted this line of work. This work was supported by a travel grant
from the Simons Foundation for the period 2026--2031.

\begin{thebibliography}{9}

\bibitem{AS03}
J.-P.~Allouche and J.~Shallit,
\emph{Automatic sequences: theory, applications, generalizations},
Cambridge University Press, Cambridge, 2003.

\bibitem{BJKP05}
V.~D.~Blondel, E.~Jeandel, P.~Koiran, and N.~Portier,
\emph{Decidable and undecidable problems about quantum automata},
SIAM J.\ Comput.\ \textbf{34} (2005), no.~6, 1464--1473.

\bibitem{CH00}
F.~Casacuberta and C.~de~la~Higuera,
\emph{Computational complexity of problems on probabilistic grammars
and transducers}, in: Grammatical Inference: Algorithms and
Applications (ICGI 2000), Lecture Notes in Computer Science 1891,
Springer, 2000, pp.~15--24.

\bibitem{GJ79}
M.~R.~Garey and D.~S.~Johnson,
\emph{Computers and intractability: a guide to the theory of
NP-completeness}, W.~H.~Freeman, San Francisco, 1979.

\bibitem{Gill24}
K.~Gill,
\emph{Probabilistic automatic complexity of finite strings},
arXiv:2402.13376 [cs.FL], 2024.

\bibitem{KH17}
B.~Kjos-Hanssen,
\emph{Superposition as memory: unlocking quantum automatic complexity},
arXiv:1703.04878 [cs.CC], 2017.

\bibitem{KHpfa}
B.~Kjos-Hanssen,
\emph{Every string has probabilistic automatic complexity at most
three}, preprint, 2026.

\bibitem{Lothaire}
M.~Lothaire,
\emph{Combinatorics on words},
Cambridge Mathematical Library, Cambridge University Press, 1997.

\bibitem{LS62}
R.~C.~Lyndon and M.-P.~Sch\"utzenberger,
\emph{The equation $a^M=b^Nc^P$ in a free group},
Michigan Math.\ J.\ \textbf{9} (1962), 289--298.

\bibitem{MC00}
C.~Moore and J.~P.~Crutchfield,
\emph{Quantum automata and quantum grammars},
Theoret.\ Comput.\ Sci.\ \textbf{237} (2000), 275--306.

\bibitem{Ren92}
J.~Renegar,
\emph{On the computational complexity and geometry of the first-order
theory of the reals, Parts I--III},
J.\ Symbolic Comput.\ \textbf{13} (1992), 255--352.

\bibitem{Sw58}
S.~\'Swierczkowski,
\emph{On a free group of rotations of the Euclidean space},
Indag.\ Math.\ \textbf{20} (1958), 376--378.

\bibitem{SW01}
J.~Shallit and M.-W.~Wang,
\emph{Automatic complexity of strings},
J.\ Autom.\ Lang.\ Comb.\ \textbf{6} (2001), no.~4, 537--554.

\bibitem{Wagon}
S.~Wagon,
\emph{The Banach--Tarski paradox},
Encyclopedia of Mathematics and its Applications 24, Cambridge
University Press, 1985.

\end{thebibliography}

\end{document}
